%% ============================================================
%%  Dispensa --- CT0540 Social Network Analysis
%%  Ca' Foscari University of Venice
%%  Prof. Fabiana Zollo
%%  A.A. 2025--2026
%% ============================================================
\documentclass[12pt,a4paper]{article}

%% ---- encoding & language ----
\usepackage[T1]{fontenc}
\usepackage[utf8]{inputenc}
\usepackage[italian]{babel}

%% ---- layout ----
\usepackage[top=2.5cm,bottom=2.5cm,left=3cm,right=3cm,headheight=15pt]{geometry}
\usepackage{setspace}
\onehalfspacing

%% ---- fonts ----
\usepackage{lmodern}
\usepackage{microtype}

%% ---- math ----
\usepackage{amsmath,amssymb,mathtools}

%% ---- code listings ----
\usepackage{listings}
\usepackage{xcolor}

\definecolor{codebg}{rgb}{0.95,0.95,0.95}
\definecolor{codecomment}{rgb}{0.4,0.4,0.4}
\definecolor{codestring}{rgb}{0.18,0.5,0.18}
\definecolor{codekeyword}{rgb}{0.0,0.35,0.7}

\lstset{
  backgroundcolor=\color{codebg},
  basicstyle=\ttfamily\small,
  commentstyle=\color{codecomment}\itshape,
  keywordstyle=\color{codekeyword}\bfseries,
  stringstyle=\color{codestring},
  breaklines=true,
  frame=single,
  framerule=0.4pt,
  xleftmargin=1em,
  xrightmargin=1em,
  aboveskip=6pt,
  belowskip=6pt,
  numbers=left,
  numberstyle=\tiny\color{gray},
  stepnumber=1,
  tabsize=2,
  showstringspaces=false
}

\lstdefinelanguage{R}{
  keywords={if,else,repeat,while,function,for,in,next,break,TRUE,FALSE,NULL,
            Inf,NaN,NA,NA_integer_,NA_real_,NA_complex_,NA_character_,
            library,require,source,return,stop,warning,message},
  sensitive=true,
  comment=[l]{\#},
  string=[b]{"'},
  morestring=[b]{'}
}

%% ---- tables ----
\usepackage{booktabs}
\usepackage{array}

%% ---- graphics & figures ----
\usepackage{graphicx}
\usepackage{float}

%% ---- hyperlinks & PDF metadata ----
\usepackage[hidelinks,
            pdftitle={Dispensa CT0540 -- Social Network Analysis},
            pdfauthor={Prof. Fabiana Zollo},
            pdfsubject={Social Network Analysis},
            pdfkeywords={reti sociali, teoria dei grafi, centralita, BERTopic, R}
            ]{hyperref}

%% ---- header / footer ----
\usepackage{fancyhdr}
\pagestyle{fancy}
\fancyhf{}
\fancyhead[L]{\small\textit{CT0540 -- Social Network Analysis}}
\fancyhead[R]{\small\textit{Ca' Foscari University of Venice}}
\fancyfoot[C]{\thepage}
\renewcommand{\headrulewidth}{0.4pt}

%% ---- section style ----
\usepackage{titlesec}
\titleformat{\section}[block]{\large\bfseries\color{blue!70!black}}{Capitolo \thesection.}{0.8em}{}[\titlerule]
\titleformat{\subsection}[block]{\normalsize\bfseries}{}{0em}{}
\titleformat{\subsubsection}[runin]{\normalsize\itshape\bfseries}{}{0em}{}[.]

%% ---- table of contents depth ----
\setcounter{tocdepth}{2}

%% ---- custom commands ----
\newcommand{\definition}[2]{%
  \vspace{4pt}%
  \noindent\textbf{Definizione (#1).}\ #2\par\vspace{4pt}%
}
\newcommand{\formula}[1]{%
  \begin{equation*}#1\end{equation*}%
}

%% ============================================================
\begin{document}

%% ---- title page ----
\begin{titlepage}
  \centering
  \vspace*{2cm}
  {\Large\bfseries Universit\`a Ca' Foscari Venezia\\[0.4em]
   Dipartimento di Scienze Ambientali, Informatica e Statistica\\[2em]}
  {\huge\bfseries\color{blue!70!black} Social Network Analysis\\[0.5em]}
  {\large CT0540 --- Dispensa del Corso\\[2em]}
  {\large Prof.ssa Fabiana Zollo\\[0.5em]}
  {\normalsize Anno Accademico 2025--2026\\[4em]}
  \hrule
  \vspace{1em}
  {\small Questa dispensa raccoglie i materiali del corso CT0540 \emph{Social Network Analysis}
  tenuto presso il Dipartimento di Scienze Ambientali, Informatica e Statistica
  dell'Universit\`a Ca' Foscari di Venezia. Il testo \`e organizzato in dieci capitoli
  che coprono progressivamente i fondamenti teorici, i metodi computazionali e le
  applicazioni empiriche dell'analisi delle reti sociali.}
  \vfill
  {\footnotesize \textit{Versione aggiornata all'Anno Accademico 2025--2026.
  Tutti i diritti riservati.}}
\end{titlepage}

\tableofcontents
\newpage

%% ============================================================
\section{Introduzione alle Reti Sociali}
%% ============================================================

\subsection{Il corso CT0540}

Questo corso si propone di introdurre gli studenti all'analisi delle reti sociali (\emph{Social Network Analysis}, SNA), una disciplina che negli ultimi due decenni ha conosciuto uno sviluppo straordinario grazie alla convergenza di tre fattori: la disponibilit\`a senza precedenti di dati digitali sulle interazioni umane, la crescente potenza di calcolo dei computer moderni, e il maturarsi di una teoria matematica ricca e rigorosa delle reti. Il corso affronta sia i fondamenti teorici --- mutuati dalla matematica, dalla sociologia e dalla fisica --- sia gli strumenti computazionali necessari per analizzare reti reali di grandi dimensioni, con particolare enfasi sull'utilizzo del linguaggio di programmazione R e del pacchetto \texttt{igraph}.

Il percorso si articola in dieci capitoli: si parte dalla definizione stessa di rete sociale e dalla sua storia intellettuale, si passa attraverso la teoria dei grafi e le sue misure fondamentali, si esaminano i meccanismi sociali che governano la formazione dei legami (omofilia, chiusura triadica, ruolo dei legami deboli), si introduce l'analisi computazionale con R, si studiano le reti di citazioni, la visualizzazione dei dati, l'analisi testuale, il topic modeling con BERTopic e, infine, le sfide metodologiche poste dai Big Data. Ogni capitolo alterna esposizione teorica, formalizzazione matematica ed esempi pratici, con l'obiettivo di rendere i concetti accessibili senza sacrificarne il rigore.

\subsection{Cosa sono le reti sociali?}

La vita sociale umana \`e intessuta di relazioni: amicizie, collaborazioni professionali, legami familiari, scambi di informazioni, transazioni economiche. Queste relazioni non esistono in isolamento ma formano strutture complesse che influenzano profondamente il comportamento individuale e collettivo. Una \emph{rete sociale} \`e precisamente la rappresentazione formale di queste strutture: un insieme di \emph{attori} (individui, organizzazioni, paesi, pagine web) collegati da \emph{relazioni} di vario tipo.

La domanda fondamentale della SNA non \`e semplicemente ``chi \`e connesso a chi?'' ma piuttosto ``in che modo la struttura delle connessioni influenza i comportamenti, le opinioni, la diffusione delle informazioni, la salute, il successo professionale?''. Questa domanda ha implicazioni profonde: se vogliamo capire come si diffonde un virus, come si diffonde una notizia falsa, come si formano le bolle informative sui social media, o come le opportunit\`a di lavoro vengono comunicate, dobbiamo capire la struttura della rete in cui questi processi avvengono.

\definition{Rete Sociale}{Una rete sociale \`e una coppia $G = (V, E)$ dove $V$ \`e un insieme di \emph{nodi} (o vertici) che rappresentano gli attori sociali, ed $E \subseteq V \times V$ \`e un insieme di \emph{archi} (o legami) che rappresentano le relazioni tra gli attori.}

\noindent Questa definizione matematica essenziale -- che approfondiremo nel capitolo dedicato alla teoria dei grafi -- cattura l'idea fondamentale: la struttura della rete \`e separabile dagli attributi individuali degli attori. Un individuo pu\`o essere ricco o povero, istruito o meno, maschio o femmina; ma la sua posizione nella rete -- quante connessioni ha, quanto \`e centrale, quanto \`e lontano dagli altri -- \`e una propriet\`a strutturale che dipende dall'intera topologia della rete, non solo dalle sue caratteristiche individuali.

\subsection{Sistemi Complessi}

Per comprendere appieno perch\'e le reti siano uno strumento cos\`i potente, occorre collocarle nel contesto pi\`u ampio dei \emph{sistemi complessi}. Aristotele scrisse nel IV secolo a.C. che ``il tutto \`e maggiore della somma delle parti'' (\emph{Metaphysica}, 1045a 8-10), e questa intuizione straordinariamente moderna cattura l'essenza dei sistemi complessi: sistemi in cui il comportamento collettivo non pu\`o essere previsto semplicemente sommando i comportamenti individuali dei componenti.

Consideriamo tre esempi emblematici. Il cervello umano contiene circa 86 miliardi di neuroni, ciascuno dei quali -- preso singolarmente -- \`e in grado di svolgere operazioni elementari di trasmissione del segnale elettrochimico. Eppure dall'interazione di questi miliardi di neuroni emergono propriet\`a straordinarie: la coscienza, la creativit\`a, la memoria, le emozioni. Nessuno di questi fenomeni \`e localizzato in un singolo neurone; sono \emph{propriet\`a emergenti} della rete neurale. Analogamente, il genoma umano conta circa 25.000 geni, che interagiscono attraverso reti di regolazione genetica di straordinaria complessit\`a: la funzione di un gene dipende spesso dall'espressione o soppressione di decine di altri geni. Infine, l'economia mondiale \`e composta da miliardi di individui e milioni di imprese che interagiscono attraverso scambi commerciali, flussi finanziari e catene di approvvigionamento: le crisi economiche sistemiche, come quella del 2008, emergono da queste interazioni in modi che nessun singolo agente poteva prevedere.

Le reti forniscono il \emph{linguaggio naturale} per descrivere i sistemi complessi. Invece di concentrarsi sulle propriet\`a individuali dei componenti, l'analisi delle reti si concentra sulla topologia delle connessioni, catturando precisamente ci\`o che fa s\`i che il tutto sia diverso dalla somma delle parti. Il concetto chiave che rende questo approccio cos\`i potente \`e quello di \textbf{universalit\`a}: reti provenienti da domini completamente diversi --- reti sociali, reti biologiche, reti tecnologiche, reti economiche --- condividono propriet\`a strutturali sorprendentemente simili. La distribuzione di grado che segue una legge di potenza (\emph{power-law}), la propriet\`a del piccolo mondo (\emph{small world}), la struttura di comunit\`a (\emph{community structure}) appaiono sistematicamente in reti tra loro molto diverse. Questa universalit\`a suggerisce che esistono principi organizzativi profondi e comuni a tutti i sistemi complessi.

\subsection{Network Science come disciplina interdisciplinare}

La \emph{network science} \`e una disciplina emergente che studia le propriet\`a e i meccanismi universali delle reti complesse. Come argomenta Albert-L\'aszl\'o Barab\'asi (Barab\'asi, 2016), la network science \`e nata come disciplina autonoma all'inizio degli anni 2000, a partire da due articoli fondamentali: il lavoro di Duncan Watts e Steven Strogatz su \emph{Nature} (1998) sui grafi piccolo-mondo, e il lavoro di Barab\'asi e Reka Albert su \emph{Science} (1999) sulle reti scale-free. Prima di questi lavori, le reti venivano studiate in modo separato e disciplinare: la sociologia aveva sviluppato la social network analysis, la matematica la teoria dei grafi, l'informatica la teoria dei grafi algoritmici, la biologia la genomica e le reti metaboliche. La network science ha unificato questi approcci, riconoscendo che le propriet\`a strutturali fondamentali delle reti emergono in tutti questi domini.

Questa interdisciplinarit\`a non \`e solo accademica: \`e una necessit\`a pratica. Per studiare, ad esempio, la diffusione di un'epidemia su una rete di contatti umani, occorre integrare la biologia dell'agente patogeno, la sociologia dei comportamenti umani, la matematica dei processi di diffusione su reti, e l'informatica per gestire dataset di milioni di individui. Nessuna di queste discipline, presa isolatamente, sarebbe sufficiente. La network science offre un quadro comune che permette a ricercatori provenienti da background diversi di comunicare e collaborare.

\subsection{Moreno e i sociogrammi (1934)}

La storia dell'analisi delle reti sociali affonda le sue radici nell'opera pionieristica di Jacob Levy Moreno (1889--1974), psichiatra e sociologo romeno naturalizzato americano. Nel 1934, Moreno pubblica \emph{Who Shall Survive?}, un'opera monumentale in cui introduce la \emph{sociometria}: la misurazione quantitativa delle relazioni interpersonali all'interno di gruppi sociali. L'idea di Moreno era rivoluzionaria per l'epoca: invece di studiare la psicologia individuale in isolamento, era necessario mappare la struttura delle relazioni interpersonali per comprendere il comportamento del gruppo.

Moreno condusse uno dei primi studi sistematici di rete sociale in una scuola di Hudson, nello stato di New York, dove chiese agli studenti di indicare accanto a quale compagno avrebbero preferito sedersi. Con queste informazioni, disegn\`o i primi \emph{sociogrammi}: diagrammi in cui gli studenti sono rappresentati come nodi e le preferenze dichiarate come archi orientati da chi preferisce verso chi \`e preferito. Questi diagrammi rivelarono immediatamente strutture sociali invisibili all'occhio nudo: nodi isolati (bambini che nessuno sceglieva e che non sceglievano nessuno), nodi molto popolari (bambini scelti da molti), preferenze reciproche (coppie che si sceglievano a vicenda), e catene di preferenze indirette.

I sociogrammi di Moreno rappresentano il \emph{primo esempio documentato di analisi delle reti sociali}. Il suo contributo \`e fondamentale per almeno tre ragioni. In primo luogo, Moreno ha dimostrato che le relazioni interpersonali hanno una struttura oggettiva che pu\`o essere misurata e rappresentata. In secondo luogo, ha mostrato che questa struttura ha conseguenze reali: i bambini isolati nel sociogramma erano spesso quelli con maggiori difficolt\`a di adattamento scolastico e sociale. In terzo luogo, ha fornito uno strumento grafico --- il sociogramma --- che rimane, nella sua essenza, la rappresentazione standard delle reti sociali ancora oggi.

\subsection{Il Club di Karate di Zachary (1977)}

Uno dei dataset pi\`u famosi e pi\`u citati nella letteratura sulla SNA \`e quello del Club di Karate di Wayne Zachary (1977). Zachary era un antropologo che studi\`o per tre anni (1970--1972) un club di karate di una universit\`a americana. Durante questo periodo, sorse un conflitto tra il maestro di karate (John A.) e il presidente del club (Hi). Il conflitto portava alla rottura del club in due fazioni: una attorno al maestro, l'altra attorno al presidente.

La cosa straordinaria \`e ci\`o che Zachary fece: mapp\`o sistematicamente tutte le interazioni sociali tra i 34 membri del club \emph{al di fuori} delle sessioni di allenamento ufficiali (conversazioni, attivit\`a condivise, amicizie al di fuori del dojo). Quando la rottura avvenne effettivamente, Zachary verific\`o che la struttura della rete -- costruita prima della rottura -- \emph{prevedeva} con straordinaria accuratezza quale fazione ciascun membro avrebbe scelto. Solo un membro su trentaquattro fu classificato in modo errato: l'individuo era strutturalmente al confine tra le due fazioni, e la sua scelta fu quindi genuinamente ambigua.

Questo risultato ha implicazioni profonde: la struttura delle relazioni sociali contiene informazioni predittive sulla dinamica futura del gruppo. Il dataset di Zachary -- 34 nodi e 78 archi -- \`e diventato il \emph{benchmark standard} per testare algoritmi di rilevamento di comunit\`a nelle reti: qualsiasi nuovo algoritmo di community detection viene tipicamente validato verificando che riesca a trovare le due comunit\`a del club di karate. La sua semplicità (dimensioni ridotte, struttura chiara) e la sua rilevanza sociologica (una vera divisione sociale documentata) lo rendono un caso di studio ideale.

\subsection{Chains of Affection (Bearman et al., 2004)}

Nel 2004, Peter Bearman, James Moody e Katherine Stovel pubblicano su \emph{American Journal of Sociology} uno studio destinato a lasciare una traccia profonda nella letteratura sulla SNA e sulla salute pubblica (Bearman, Moody \& Stovel, 2004). Lo studio mappc\`o la rete delle relazioni romantiche e sessuali tra gli studenti di una scuola superiore americana di medie dimensioni, che gli autori chiamarono Jefferson High, situata in una citt\`a di circa 30.000 abitanti nel Midwest americano.

I dati furono raccolti attraverso interviste approfondite con 573 studenti, che dichiararono i loro partner negli ultimi diciotto mesi. Il risultato fu sorprendente: invece delle dense strutture a cluster che ci si sarebbe aspettati (molte coppie separate, con pochi ponti), la rete presentava una struttura a \emph{spanning tree} -- un lungo percorso che attraversava l'intera rete, toccando quasi met\`a dei partecipanti. Questa struttura, che gli autori chiamarono ``catena di affetti'' (\emph{chain of affection}), aveva implicazioni drammatiche per la trasmissione delle malattie sessualmente trasmissibili: una persona infetta a un'estremit\`a della catena poteva potenzialmente trasmettere l'infezione a centinaia di altre persone attraverso percorsi relativamente brevi, anche se non aveva mai avuto contatti diretti con loro.

Lo studio di Bearman e colleghi \`e un esempio paradigmatico di come la struttura della rete -- non solo i comportamenti individuali -- determini i rischi collettivi. La conclusione pratica \`e che gli interventi di salute pubblica devono tenere conto della topologia della rete di trasmissione, non solo del numero di partner individuali.

\subsection{Computational Social Science (Lazer et al., 2009)}

Nel febbraio 2009, un gruppo di ricercatori guidato da David Lazer pubblica su \emph{Science} un articolo intitolato ``Computational Social Science'' (Lazer et al., 2009), che propone formalmente l'esistenza di una nuova disciplina scientifica nata dall'incontro tra le scienze sociali e l'informatica. L'argomento centrale \`e che la proliferazione di \emph{tracce digitali} -- email, telefonate, dati GPS, post sui social media, transazioni economiche -- stia aprendo possibilit\`a senza precedenti per lo studio del comportamento umano a scale e con risoluzioni temporali mai viste prima.

Le scienze sociali tradizionali si basavano su metodi come sondaggi, esperimenti di laboratorio e osservazione etnografica, ciascuno dei quali ha limitazioni significative: i sondaggi sono costosi, rari e soggetti a bias di risposta; gli esperimenti di laboratorio sono artificiali e coinvolgono campioni piccoli; l'etnografia \`e profonda ma non scalabile. Le tracce digitali, al contrario, sono generate naturalmente dai comportamenti quotidiani di milioni di persone, sono continue nel tempo, e hanno una granularit\`a straordinaria: possiamo sapere non solo che una persona \`e andata al lavoro, ma esattamente quale percorso ha preso, con chi ha interagito, quali messaggi ha scritto.

Questo cambiamento ha trasformato profondamente la ricerca sulle reti sociali. Dataset come quelli dei social media (Twitter, Facebook, Instagram), delle piattaforme di messaggistica (WhatsApp, Telegram), delle reti di collaborazione scientifica (co-authorship networks) o delle reti di citazioni accademiche hanno permesso di studiare fenomeni sociali a scale che sarebbero state impensabili anche solo vent'anni fa. Allo stesso tempo, questa abbondanza di dati porta con s\`e nuove sfide metodologiche ed etiche, che verranno discusse in dettaglio nel capitolo dedicato ai Big Data.

\subsection{Il numero di Dunbar ($\sim 150$)}

Quante relazioni sociali stabili pu\`o mantenere un essere umano? Questa domanda, che pu\`o sembrare banale, ha una risposta sorprendentemente precisa grazie alla ricerca dell'antropologo Robin Dunbar. Negli anni '90, Dunbar studi\`o la dimensione del neocorteccia --- la parte evolutivamente pi\`u recente del cervello, associata alle funzioni cognitive superiori --- in diverse specie di primati, e la mise in relazione con la dimensione media dei loro gruppi sociali. Trov\`o una correlazione sorprendentemente robusta: i primati con neocorteccia pi\`u grande vivono in gruppi sociali pi\`u grandi. Estrapolando questa relazione alla dimensione del neocorteccia umana, Dunbar calcol\`o che gli esseri umani dovrebbero essere in grado di mantenere relazioni sociali stabili con circa 150 individui (Dunbar, 1992).

Questo numero, noto come \emph{numero di Dunbar}, ha trovato conferme empiriche in ambiti molto diversi: la dimensione media dei villaggi neolitici, dei reggimenti militari in tutte le epoche storiche, delle unit\`a operative delle aziende. La spiegazione evolutiva \`e che mantenere una relazione sociale stabile richiede un investimento cognitivo significativo --- ricordare chi \`e la persona, quale \`e la storia della relazione, quali sono le sue preferenze e i suoi stati emotivi --- e che la capacit\`a cognitiva disponibile pone un limite naturale.

La verifica del numero di Dunbar nell'era dei social media \`e stata condotta da Goncalves e colleghi (Goncalves, Perra \& Flammini, 2011), che hanno analizzato i dati di 1,7 milioni di utenti Twitter. Nonostante molti utenti abbiano migliaia di follower, gli autori hanno trovato che ciascun utente interagisce attivamente --- cio\`e scambia messaggi reciproci in modo regolare --- con un numero di persone compreso tra 100 e 200, con una media vicina a 150. Questo suggerisce che i social media non hanno eliminato il limite cognitivo del numero di Dunbar: hanno semplicemente spostato il \emph{tipo} di relazioni che possiamo avere, permettendo relazioni pi\`u deboli e superficiali con un numero molto maggiore di persone, ma senza aumentare il numero di relazioni davvero significative.

\subsection{Information Networks}

Le reti di informazione sono reti in cui i nodi rappresentano risorse informative piuttosto che attori sociali, e i legami rappresentano riferimenti, citazioni, o link tra queste risorse. Il World Wide Web \`e forse la rete di informazione pi\`u grande mai creata: miliardi di pagine web collegate da hyperlink orientati (da una pagina verso le pagine a cui rimanda). Le reti di citazioni accademiche sono un altro esempio fondamentale: gli articoli scientifici sono i nodi, e i riferimenti bibliografici sono gli archi orientati che puntano dalla pubblicazione che cita a quella citata.

Lo studio della blogosfera americana prima delle elezioni presidenziali del 2004 condotto da Adamic e Glance (Adamic \& Glance, 2005) \`e un lavoro seminale nell'analisi delle reti di informazione politica. Gli autori raccolsero i dati di oltre 1.400 blog politici americani, classificandoli come liberali o conservatori sulla base dei loro contenuti, e mapparono i link reciproci tra di essi. Il risultato fu una polarizzazione politica netta e sorprendente: i blog liberali linkavano principalmente altri blog liberali, e i blog conservatori linkavano principalmente altri blog conservatori. Le due comunit\`a erano quasi completamente separate, con pochissimi ponti tra loro. Questo fenomeno, che Adamic e Glance definirono \emph{echo chamber} --- camera d'eco --- si \`e rivelato un precursore di dinamiche che oggi conosciamo bene sui social media.

Il concetto di echo chamber nel dibattito politico su Facebook \`e stato studiato empiricamente da Del Vicario, Zollo e colleghi (Del Vicario, Zollo et al., 2017), in uno studio sul dibattito attorno alla Brexit sulla piattaforma Facebook. Lo studio ha mostrato che gli utenti tendevano a consumare e a condividere informazioni allineate con le proprie opinioni preesistenti, formando comunit\`a di utenti che si esponevano prevalentemente a prospettive omogenee. Questo meccanismo -- che combina la selezione attiva delle fonti da parte degli utenti con la personalizzazione algoritmica dei feed --- amplifica la polarizzazione politica e rende pi\`u difficile il dialogo tra posizioni diverse.

\subsection{Tipi di reti}

Le reti possono essere classificate in diversi modi a seconda della natura dei nodi e degli archi. Una prima distinzione fondamentale riguarda la \emph{direzione} degli archi: una rete \`e \emph{orientata} (o \emph{diretta}) quando gli archi hanno una direzione (da $u$ a $v$ non \`e lo stesso che da $v$ a $u$), e \emph{non orientata} quando i legami sono simmetrici. Le amicizie su Facebook sono tipicamente non orientate (se A \`e amico di B, allora B \`e amico di A), mentre i ``follower'' su Twitter sono orientati (A pu\`o seguire B senza che B segua A).

Una seconda distinzione riguarda la presenza di \emph{pesi} sugli archi: una rete \`e \emph{pesata} quando i legami hanno intensit\`a diverse (ad esempio, due persone che si telefonano cento volte al mese hanno un legame pi\`u forte di due persone che si telefonano una volta). Una rete \`e \emph{non pesata} (o binaria) quando i legami indicano semplicemente la presenza o assenza di una relazione.

Oltre alle reti sociali in senso stretto, \`e utile distinguere le \emph{reti di informazione} (come il Web e le reti di citazioni), le \emph{reti tecnologiche} (come Internet, le reti elettriche, le reti di trasporto), le \emph{reti biologiche} (reti di proteine, reti metaboliche, reti neurali) e le \emph{reti economiche} (reti di scambi commerciali, reti interbancarie). Come vedremo nel corso, queste reti condividono molte propriet\`a strutturali fondamentali nonostante la loro diversit\`a apparente.

\subsection{Testi di riferimento}

I testi fondamentali a cui questo corso fa riferimento sono i seguenti. Per i fondamenti della network science: Barab\'asi, A.-L. (2016), \emph{Network Science}, Cambridge University Press (disponibile gratuitamente online). Per la SNA classica: Easley, D. \& Kleinberg, J. (2010), \emph{Networks, Crowds, and Markets}, Cambridge University Press. Per l'implementazione computazionale: Ognyanova, K. (materiali del corso, disponibili online). Per il topic modeling: Grootendorst, M. (2022), ``BERTopic: Neural topic modeling with a class-based TF-IDF procedure'', \emph{arXiv:2203.05794}. Per i Big Data: Boyd, D. \& Crawford, K. (2012), ``Critical Questions for Big Data'', \emph{Information, Communication \& Society}.

\subsection{Ricerca e progetti del gruppo}

Il gruppo di ricerca della Prof.ssa Zollo si occupa di analisi computazionale dei fenomeni di disinformazione, polarizzazione e dinamiche di opinione sui social media. Tra i contributi pi\`u rilevanti: Del Vicario et al. (2016) hanno studiato la diffusione di notizie vere e false su Facebook, mostrando che le \emph{echo chambers} favoriscono la propagazione della disinformazione; Zollo et al. (2017) hanno analizzato come gli utenti reagiscono alle smentite (\emph{debunking}) e come le reazioni differiscano tra comunit\`a con credenze diverse; Pi\~nero et al. (2023) hanno applicato tecniche di NLP avanzate per la rilevazione automatica della disinformazione multilingue. Questi contributi si inseriscono nel contesto pi\`u ampio della Computational Social Science, che vedremo in dettaglio nel capitolo sui Big Data.


%% ============================================================
\section{Teoria dei Grafi}
%% ============================================================

\subsection{I ponti di K\"{o}nigsberg}

La teoria dei grafi nasce ufficialmente nel 1736 con un articolo di Leonhard Euler intitolato ``Solutio problematis ad geometriam situs pertinentis'' (\emph{Commentarii Academiae Scientiarum Imperialis Petropolitanae}, 1736). Il problema che Euler risolse --- o meglio, dimostr\`o essere irrisolvibile --- era quello dei sette ponti di K\"onigsberg: \`e possibile percorrere tutti e sette i ponti che attraversano il fiume Pregel nella citt\`a di K\"onigsberg (oggi Kaliningrad) passando su ciascuno esattamente una volta? I cittadini di K\"onigsberg ci avevano provato per anni senza riuscirci, e Euler dimostr\`o formalmente che era impossibile.

Il contributo geniale di Euler fu di astrazione: invece di ragionare sulla geografia fisica della citt\`a, rappresent\`o il problema come un \emph{grafo} in cui le masse di terra (due isole e le due rive del fiume) sono nodi e i ponti sono archi. Il problema diventa cos\`i: \`e possibile trovare un percorso nel grafo che attraversi ogni arco esattamente una volta (oggi chiamato \emph{cammino euleriano})? Euler dimostr\`o che un tale percorso esiste se e solo se il grafo ha esattamente zero o due nodi con grado dispari. Nel caso di K\"onigsberg, tutti e quattro i nodi avevano grado dispari (tre ponti ciascuno), quindi il percorso era impossibile.

Questa soluzione ha un significato che va ben oltre il problema specifico: Euler ha mostrato che alcune propriet\`a della rete (la distribuzione dei gradi) determinano fenomeni macroscopici (la possibilit\`a o meno di un percorso euleriano) indipendentemente da tutti gli altri dettagli della struttura. Questo principio --- le propriet\`a topologiche della rete hanno conseguenze globali --- \`e al cuore di tutta la moderna teoria delle reti.

\subsection{Definizioni fondamentali}

Prima di procedere con misure pi\`u avanzate, \`e necessario stabilire con precisione il vocabolario di base della teoria dei grafi. Queste definizioni non sono mere formalit\`a accademiche: ciascuna cattura un aspetto distinto della struttura di una rete, e la precisione terminologica \`e essenziale per evitare ambiguit\`a nell'interpretazione dei risultati.

\definition{Grafo}{Un grafo $G = (V, E)$ \`e una coppia ordinata composta da un insieme $V$ di \emph{vertici} (o \emph{nodi}) e un insieme $E$ di \emph{archi} (o \emph{spigoli}), dove ogni arco \`e una coppia non ordinata $\{u, v\}$ di vertici distinti. Indichiamo con $n = |V|$ il numero di nodi e con $m = |E|$ il numero di archi.}

\noindent La distinzione tra grafi orientati e non orientati \`e cruciale in molte applicazioni. In un grafo non orientato, il legame tra $u$ e $v$ \`e lo stesso del legame tra $v$ e $u$: le amicizie, le collaborazioni, la prossimit\`a fisica sono esempi di relazioni simmetriche. In un grafo orientato (\emph{digrafo}), invece, ogni arco ha una direzione: il ``seguire'' su Twitter, il ``citare'' in un articolo scientifico, il ``trasferire denaro'' in una transazione finanziaria sono esempi di relazioni asimmetriche.

\definition{Grafo Orientato}{Un grafo orientato (o \emph{digrafo}) \`e una coppia $G = (V, E)$ dove gli archi in $E$ sono coppie \emph{ordinate} $(u, v)$ con $u \neq v$. L'arco $(u, v)$ va da $u$ (sorgente) a $v$ (destinazione), ed \`e distinto dall'arco $(v, u)$.}

\definition{Grafo Pesato}{Un grafo pesato \`e un grafo $G = (V, E, w)$ in cui ad ogni arco $(u, v) \in E$ \`e associato un \emph{peso} $w(u,v) \in \mathbb{R}$, che pu\`o rappresentare la forza del legame, la frequenza di interazione, la distanza fisica, o qualsiasi altra quantit\`a rilevante per il fenomeno studiato.}

\definition{Adiacenza}{Due nodi $u$ e $v$ si dicono \emph{adiacenti} (o \emph{vicini}) se esiste un arco che li connette. L'insieme di tutti i vicini di un nodo $v$ si chiama \emph{vicinato} di $v$ e si denota $N(v) = \{u \in V : \{u,v\} \in E\}$.}

\subsection{Grado e distribuzione di grado}

Prima di addentrarci nelle misure pi\`u sofisticate, \`e necessario partire dal concetto pi\`u semplice: il \textbf{grado} di un nodo. Immaginate di trovarvi a una festa: il vostro grado corrisponderebbe al numero di persone con cui avete scambiato almeno un saluto. In termini formali, il grado $k_i$ di un nodo $i$ \`e il numero di archi che collegano $i$ agli altri nodi della rete.

\definition{Grado}{Il \emph{grado} $k_i$ di un nodo $i$ in un grafo non orientato \`e definito come il numero di archi incidenti su $i$:
$k_i = |\{j \in V : \{i,j\} \in E\}| = |N(i)|$.}

\noindent Una relazione fondamentale lega la somma dei gradi al numero di archi. Poich\'e ogni arco contribuisce esattamente due al conteggio totale dei gradi (uno per ciascuno dei due nodi estremi), abbiamo:

\formula{\sum_{i \in V} k_i = 2m}

\noindent dove $m = |E|$ \`e il numero totale di archi. Questa identit\`a, nota come \emph{lemma delle strette di mano}, ha un'interpretazione intuitiva immediata: ogni legame contribuisce al grado di entrambi i nodi che connette, quindi la somma dei gradi conta ogni legame due volte. Da questa relazione si ricava immediatamente che il grado medio di un nodo \`e:

\formula{\langle k \rangle = \frac{1}{n} \sum_{i \in V} k_i = \frac{2m}{n}}

\noindent dove $n$ \`e il numero totale di nodi. Questa formula ci dice che il grado medio dipende dal rapporto tra il numero di archi e il numero di nodi: una rete densa (con molti archi rispetto ai nodi) ha grado medio alto, una rete sparsa ha grado medio basso.

Nei grafi orientati, si distinguono due tipi di grado. Il \emph{grado entrante} (\emph{in-degree}) $k_i^{\text{in}}$ di un nodo $i$ \`e il numero di archi che arrivano su $i$, cio\`e il numero di nodi che puntano verso $i$. Il \emph{grado uscente} (\emph{out-degree}) $k_i^{\text{out}}$ \`e il numero di archi che partono da $i$, cio\`e il numero di nodi verso cui $i$ punta. Per un grafo orientato vale:

\formula{\sum_{i \in V} k_i^{\text{in}} = \sum_{i \in V} k_i^{\text{out}} = m}

\noindent Questa identit\`a riflette il fatto che ogni arco orientato ha esattamente una sorgente (contribuisce a un grado uscente) e una destinazione (contribuisce a un grado entrante).

La \emph{distribuzione di grado} $P(k)$ di una rete descrive la probabilit\`a che un nodo scelto a caso abbia grado $k$. \`E la quantit\`a statistica pi\`u fondamentale per caratterizzare la struttura di una rete. Per le reti reali di grande interesse (Internet, Web, reti di collaborazione scientifica, reti biologiche), la distribuzione di grado segue approssimativamente una \emph{legge di potenza}:

\formula{P(k) \sim k^{-\gamma}}

\noindent dove $\gamma$ \`e un esponente che tipicamente cade nell'intervallo $2 < \gamma < 3$ per le reti reali. Le reti con distribuzione di grado a legge di potenza vengono dette \emph{scale-free} (prive di scala) perch\'e non hanno una scala caratteristica del grado: non esiste un ``grado tipico'' come avviene invece in una distribuzione gaussiana. Le reti scale-free sono caratterizzate da una piccola minoranza di nodi con grado elevatissimo (gli \emph{hub}) e una grande maggioranza di nodi con grado basso. Questo tipo di distribuzione ha implicazioni profonde per la robustezza delle reti (le reti scale-free sono robuste a guasti casuali ma vulnerabili ad attacchi mirati agli hub) e per la diffusione di epidemie e informazioni.

\subsection{Matrice di adiacenza}

Per poter lavorare computazionalmente con un grafo, \`e necessario rappresentarlo in modo che un computer possa manipolarlo. La rappresentazione pi\`u naturale e pi\`u usata nell'algebra lineare \`e la \emph{matrice di adiacenza}.

\definition{Matrice di Adiacenza}{La matrice di adiacenza di un grafo $G = (V, E)$ con $n$ nodi \`e una matrice $A \in \{0,1\}^{n \times n}$ definita da:
$A_{ij} = 1$ se $\{i,j\} \in E$ (esiste un arco tra $i$ e $j$), e $A_{ij} = 0$ altrimenti.}

\noindent Per un grafo non orientato, la matrice di adiacenza \`e simmetrica: $A_{ij} = A_{ji}$ per ogni coppia $(i,j)$. Per un grafo orientato, $A_{ij} = 1$ indica l'esistenza dell'arco orientato da $i$ verso $j$, e la matrice non \`e necessariamente simmetrica.

Una propriet\`a elegante della matrice di adiacenza \`e che le sue potenze hanno un'interpretazione reteologica precisa. La potenza $A^k$ \`e la matrice il cui elemento $(i,j)$ conta il numero di cammini di lunghezza esattamente $k$ da $i$ a $j$. In particolare, $(A^2)_{ij} = \sum_{\ell} A_{i\ell} A_{\ell j}$ conta il numero di nodi che sono vicini sia di $i$ che di $j$ --- cio\`e il numero di vicini comuni tra $i$ e $j$.

La matrice di adiacenza permette anche di calcolare facilmente il grado di ogni nodo: il grado di $i$ \`e semplicemente la somma degli elementi della riga $i$ (o della colonna $i$, per grafi non orientati):

\formula{k_i = \sum_{j=1}^{n} A_{ij}}

\noindent Tuttavia, per reti molto grandi, la matrice di adiacenza diventa impraticabile: una rete con $n$ nodi richiede $n^2$ celle di memoria, e per $n = 10^6$ questo corrisponde a $10^{12}$ celle --- impensabile con hardware standard. Per questo motivo, le reti reali grandi vengono quasi sempre rappresentate tramite \emph{liste di adiacenza}.

\subsection{Liste di adiacenza}

Una lista di adiacenza \`e una struttura dati in cui, per ogni nodo $i$, si memorizza esplicitamente solo l'insieme dei suoi vicini $N(i)$. Il costo di memoria totale \`e $O(n + m)$, dove $n$ \`e il numero di nodi e $m$ il numero di archi. Per reti sparse (dove $m \ll n^2$), questo \`e enormemente pi\`u efficiente della matrice di adiacenza.

La scelta tra matrice di adiacenza e lista di adiacenza dipende dal tipo di operazioni che dobbiamo eseguire. La matrice di adiacenza permette di verificare in $O(1)$ se esiste un arco tra due nodi specifici, ma richiede $O(n^2)$ di spazio. La lista di adiacenza richiede $O(\deg(i))$ per verificare se esiste un arco da $i$ a $j$, ma usa solo $O(n + m)$ di spazio e \`e molto pi\`u efficiente per operazioni che devono scorrere tutti i vicini di un nodo (come il calcolo del grado, la BFS, la DFS).

\subsection{Reti sparse e dense}

Una rete si dice \emph{sparsa} quando il numero di archi $m$ \`e molto inferiore al numero massimo possibile di archi $\binom{n}{2} = \frac{n(n-1)}{2}$ per grafi non orientati. La \emph{densit\`a} di un grafo \`e definita come:

\formula{\delta = \frac{2m}{n(n-1)}}

\noindent Per un grafo completo (dove ogni coppia di nodi \`e connessa), $\delta = 1$. Per le reti reali, la densit\`a \`e tipicamente molto bassa: la rete dei paper scientifici di un campo ha densit\`a dell'ordine di $10^{-4}$ o inferiore, il Web ha densit\`a dell'ordine di $10^{-9}$ o inferiore. Le reti sparse hanno propriet\`a strutturali molto diverse dalle reti dense, e gli algoritmi efficienti per le reti reali grandi devono sfruttare la sparsit\`a.

\subsection{Connettivit\`a e componenti connesse}

Un aspetto fondamentale della struttura di una rete \`e la sua \emph{connettivit\`a}: \`e possibile raggiungere ogni nodo partendo da ogni altro nodo?

\definition{Componente Connessa}{Una \emph{componente connessa} di un grafo non orientato \`e un insieme massimale di nodi tali che esiste un cammino tra ogni coppia di nodi dell'insieme.}

\noindent Per i grafi orientati, si distinguono le \emph{componenti fortemente connesse} (in cui si pu\`o raggiungere ogni nodo da ogni altro nodo seguendo gli archi nella loro direzione) e le \emph{componenti debolmente connesse} (in cui si pu\`o raggiungere ogni nodo da ogni altro ignorando la direzione degli archi).

\subsection{La componente gigante}

Nelle reti reali di grandi dimensioni, si osserva tipicamente la presenza di una \emph{componente gigante}: una singola componente connessa che contiene una frazione macroscopica (spesso il 50--90\%) di tutti i nodi della rete, mentre le restanti componenti sono piccole (di dimensione $O(\log n)$ o costante). La presenza della componente gigante \`e un fenomeno di transizione di fase: nel modello di Erd\H{o}s-R\'enyi (rete casuale in cui ogni coppia di nodi \`e connessa indipendentemente con probabilit\`a $p$), la componente gigante appare bruscamente quando il grado medio supera 1, cio\`e quando $\langle k \rangle = pn > 1$.

Questa transizione ha un'interpretazione intuitiva: se ogni nodo ha in media meno di un vicino, la rete \`e frammentata in molte piccole isole disconnesse. Ma non appena il grado medio supera 1, queste isole si fondono in una componente gigante che attraversa l'intera rete. La stessa dinamica si osserva nella diffusione di epidemie: se il numero di riproduzione di base $R_0 > 1$, l'epidemia si diffonde nell'intera popolazione; se $R_0 < 1$, si spegne spontaneamente.

\subsection{Cammini, distanza e diametro}

\definition{Cammino}{Un \emph{cammino} in un grafo $G = (V, E)$ \`e una sequenza di nodi $v_0, v_1, v_2, \ldots, v_\ell$ tale che per ogni $i \in \{0, 1, \ldots, \ell-1\}$ esiste un arco $\{v_i, v_{i+1}\} \in E$. La \emph{lunghezza} del cammino \`e il numero di archi attraversati, cio\`e $\ell$.}

\definition{Distanza Geodesica}{La \emph{distanza geodesica} $d(u, v)$ tra due nodi $u$ e $v$ \`e la lunghezza del cammino pi\`u breve (o \emph{geodetico}) che collega $u$ a $v$. Se non esiste alcun cammino, si pone $d(u,v) = \infty$.}

\noindent La lunghezza media del percorso (\emph{average path length}) di una rete \`e:

\formula{\bar{L} = \frac{1}{n(n-1)} \sum_{i \neq j} d(i,j)}

\noindent Questa quantit\`a misura la ``distanza tipica'' tra due nodi scelti a caso nella rete: \`e una misura globale della coesione della rete. Il \emph{diametro} $D$ di una rete \`e la massima distanza geodesica tra qualsiasi coppia di nodi:

\formula{D = \max_{i,j \in V} d(i,j)}

\noindent Il diametro misura la ``larghezza'' della rete: il caso peggiore. Una rete con diametro piccolo \`e ``compatta'': da qualsiasi nodo si riesce a raggiungere qualsiasi altro in pochi passi.

\subsection{Il Piccolo Mondo: Milgram, Watts e Strogatz}

Il fenomeno del ``piccolo mondo'' (small world) \`e uno dei risultati pi\`u sorprendenti e controintuitivi dell'analisi delle reti. L'intuizione comune porterebbe a pensare che in una rete di 7 miliardi di persone, la distanza media tra due individui qualsiasi sia enorme. Stanley Milgram, psicologo sociale di Harvard, test\`o sperimentalmente questa intuizione nel 1967. Chiese a persone in Nebraska e Kansas di far arrivare una lettera a un bersaglio a Boston, passando la lettera a un conoscente che si pensava fosse pi\`u vicino al bersaglio. Il risultato fu straordinario: le catene che arrivarono a destinazione avevano in media circa 5--6 passaggi, da cui l'espressione ``sei gradi di separazione''.

Molti decenni dopo, Duncan Watts e Steven Strogatz (Watts \& Strogatz, 1998) formalizzarono matematicamente il fenomeno del piccolo mondo. Essi dimostrarono che molte reti reali soddisfano contemporaneamente due propriet\`a che a prima vista sembrano incompatibili: una \emph{distanza media corta} (come i grafi casuali di Erd\H{o}s-R\'enyi) e un \emph{elevato coefficiente di clustering} (come i grafi regolari, dove i vicini dei vicini tendono a essere vicini). I grafi casuali hanno distanza media corta ma clustering basso; i grafi regolari hanno clustering alto ma distanza media lunga. Le reti reali --- e il modello di Watts-Strogatz --- hanno entrambe le propriet\`a.

\definition{Coefficiente di Clustering Globale}{Il \emph{coefficiente di clustering} $C$ di un grafo misura la proporzione di triangoli chiusi rispetto al numero di ``triadi'' aperte (cammini di lunghezza 2). Formalmente:
$C = \frac{3 \times \text{numero di triangoli}}{\text{numero di triadi aperte}}$}

\noindent Un'interpretazione pi\`u intuitiva si ottiene dal coefficiente di clustering locale. Il coefficiente di clustering locale $C_i$ di un nodo $i$ misura la proporzione di coppie di vicini di $i$ che sono connesse tra loro:

\formula{C_i = \frac{|\{(u,v) : u,v \in N(i),\, \{u,v\} \in E\}|}{k_i(k_i-1)/2}}

\noindent Il numeratore conta il numero di archi effettivamente presenti tra i vicini di $i$; il denominatore conta il numero massimo possibile di tali archi (tutte le coppie di vicini). Se $C_i = 1$, tutti i vicini di $i$ sono connessi tra loro (il vicinato di $i$ forma un grafo completo). Se $C_i = 0$, nessuna coppia di vicini di $i$ \`e connessa.

\subsection{Reti scale-free e il modello di Barab\'asi-Albert}

Come si spiegano le reti scale-free con distribuzione di grado a legge di potenza? Barab\'asi e Albert (Barab\'asi \& Albert, 1999) proposero un meccanismo generativo elegante basato su due ingredienti: la \emph{crescita} (la rete cresce nel tempo, con nuovi nodi che si aggiungono continuamente) e il \emph{preferential attachment} (i nuovi nodi tendono a connettersi preferenzialmente ai nodi gi\`a molto connessi --- ``i ricchi diventano pi\`u ricchi'').

Pi\`u precisamente, nel modello di Barab\'asi-Albert (BA), ad ogni passo temporale viene aggiunto un nuovo nodo che si connette a $m$ nodi gi\`a presenti nella rete con una probabilit\`a proporzionale al loro grado attuale:

\formula{\Pi(k_i) = \frac{k_i}{\sum_j k_j}}

\noindent dove $\Pi(k_i)$ \`e la probabilit\`a che il nuovo nodo si connetta al nodo $i$, e $k_i$ \`e il grado attuale di $i$. Questo meccanismo genera, nel limite di reti grandi, una distribuzione di grado a legge di potenza $P(k) \sim k^{-3}$. L'intuizione \`e chiara: i nodi che arrivano prima nella rete hanno il tempo di accumulare pi\`u connessioni, e la loro maggiore connettivit\`a li rende ancora pi\`u attrattivi per i nuovi arrivati. Il risultato \`e una distribuzione estremamente eterogenea con pochi hub dominanti.

Il preferential attachment \`e un meccanismo plausibile per molte reti reali: le pagine web pi\`u famose ricevono pi\`u link (perch\'e gli utenti le conoscono e le citano); gli articoli scientifici gi\`a molto citati ricevono nuove citazioni pi\`u facilmente (perch\'e appaiono ai primi posti nelle ricerche bibliografiche); le persone gi\`a famose sui social media acquisiscono nuovi follower pi\`u velocemente. Questo meccanismo di ``feedback positivo'' \`e ubiquo nei sistemi complessi e produce la caratteristica distribuzione a legge di potenza osservata empiricamente.

